\documentclass[11pt]{article}
\usepackage{ochamath}
\subjectcolor{2F5DA8}
\setsheet{線形代数}{行列式　演習}{https://university-math-crj.pages.dev/linear-algebra/determinants/}
\begin{document}
\makesheethead{解答}

\problem[2×2の計算]
\begin{ans}
$\det A=3\times2-1\times2=4$。0ではないので正則。単位正方形の像の面積は4。
\end{ans}
\point{符号付きの値と面積の絶対値を区別する。}

\problem[三角行列]
\begin{ans}
三角行列なので $\det T=2\times(-1)\times4=-8$。体積倍率は $|-8|=8$。符号は向きの反転を表す。
\end{ans}
\point{三角行列の非対角成分は行列式の積に入らない。}
\newpage

\problem[行操作]
\begin{ans}
交換で $-4$、第1行の3倍で $-12$。実際、行列は $\begin{pmatrix}6&6\\3&1\end{pmatrix}$ で、$6-18=-12$。
\end{ans}
\point{全行を3倍する場合とは倍率が異なる。}

\problem[面積と形]
\begin{ans}
$\det D=1$ で面積を保つ。しかし $(1,0)^{\mathsf T}$ は $(2,0)^{\mathsf T}$ に移り、長さは1から2になる。長さを保つという主張は誤り。
\end{ans}
\point{1本の反例で、すべてのベクトルに対する主張を否定できる。}
\end{document}
